% LREC 2026 paper draft aligned with current repository outputs
\documentclass[10pt, a4paper]{article}

\usepackage{booktabs}
\usepackage{array}
\usepackage{multirow}
\usepackage{makecell}
\usepackage{graphicx}
\usepackage[T1]{fontenc}
\usepackage{lmodern}
\usepackage{amsmath,amssymb}
\usepackage[review]{lrec2026}

\title{What Signals Really Matter for Misinformation Tasks? \\
Evaluation-First Benchmarks for Fake-News Detection and Virality Prediction}

\name{Francesco Paolo Savatteri, Chahan Vidal-Gor\`ene, Florian Cafiero}

\address{Ecole nationale des chartes - PSL\\
65 rue de Richelieu, 72002, Paris, France \\
francesco.savatteri@chartes.psl.edu, chahan.vidalgorene@chartes.psl.eu, florian.cafiero@chartes.psl.eu}

\abstract{
We present an evaluation-first study of two operational misinformation tasks: (i) fake-news detection and (ii) virality prediction. Using EVONS and FakeNewsNet (Politifact), we compare frozen text-embedding pipelines (RoBERTa and Mistral) with lightweight downstream classifiers (MLP, logistic regression, random forest, and XGBoost when available). We report consolidated 10-split results from the repository artifacts and complement descriptive metrics with paired inferential tests (bootstrap confidence intervals, sign-flip permutation tests, Holm--Bonferroni correction, Cliff's $\delta$). The results show a robust task-difficulty gap: fake-news detection is near ceiling on EVONS and strong on FakeNewsNet, whereas virality prediction is more sensitive to label construction and class imbalance. On EVONS virality (95th-percentile label), only \texttt{gating\_mistral} attains substantial F1 among tested variants. On FakeNewsNet virality (median split), several models are close, with \texttt{xgboost\_bert} leading by mean F1 but without corrected significance over the strongest alternatives. These findings support transparent, statistically controlled reporting before architecture-centric claims.
\\ \newline \Keywords{Misinformation detection, Virality prediction, Evaluation protocol, Statistical significance, EVONS, FakeNewsNet, Reproducibility}
}

\begin{document}
\maketitleabstract

\section{Introduction}
The rapid spread of misleading online content has motivated two complementary goals: detecting misinformation and anticipating diffusion potential (``virality'') in time for intervention \citep{Vosoughi2018,Lazer2018,Shu2017}. Although recent literature emphasizes multimodal and graph-based systems, practical deployment still hinges on foundational design choices: (i) how labels are operationalized, (ii) which signals are realistically available, and (iii) whether reported gains survive statistical controls.

This work takes an \emph{evaluation-first} position. We do not introduce a new architecture. Instead, we align and report the current repository pipeline and outputs in a coherent benchmark narrative, emphasizing reproducibility, comparability, and inferential validity.

\section{Tasks and Data}
\label{sec:tasks-data}

\paragraph{Task T1: Fake-news detection.}
Binary veracity classification (fake vs. real).

\paragraph{Task T2: Virality prediction.}
Binary classification of viral vs. non-viral items.

\paragraph{EVONS.}
EVONS contains article-level text (title/description), source metadata, veracity labels, and aggregate Facebook engagement. For virality, positives are defined by the upper tail of engagement (95th percentile), producing strong class imbalance.

\paragraph{FakeNewsNet (Politifact).}
The repository uses preprocessed Politifact instances. For virality, labels are induced by a median split on total likes, yielding approximately balanced classes.

\section{Methods}
\label{sec:methods}

\paragraph{Representations.}
Texts are embedded upstream and reused during model training. Two embedding backends are compared: RoBERTa-base (768 dimensions) and Mistral embeddings (1024 dimensions).

\paragraph{Model families.}
Across tasks, we evaluate lightweight downstream classifiers: MLP, logistic regression, random forest, and XGBoost (when available). On EVONS virality, additional source-aware MLP variants are included (source embedding, average engagement feature, and gated fusion).

\paragraph{Evaluation protocol.}
All reported runs use repeated stratified hold-out with \texttt{StratifiedShuffleSplit} (10 splits, 80/20 train-validation), fixed random seed, and standard metrics (Accuracy, F1, Precision, Recall, ROC-AUC). In imbalanced settings, interpretation prioritizes F1/Recall and ROC-AUC over raw Accuracy.

\paragraph{Statistical controls.}
Model comparisons are tested pairwise with sign-flip permutation tests and bootstrap confidence intervals (10,000 resamples). Multiple testing is controlled with Holm--Bonferroni; effect sizes are reported with Cliff's $\delta$.

\section{Results}
\label{sec:results}

Table~\ref{tab:readable-full-results} reports consolidated means across 10 splits for all task/dataset/model settings. Table~\ref{tab:stats-top-pairs} reports top F1 pairwise comparisons (best model vs. strongest alternatives) with inferential controls.

% Auto-generated by tools/generate\_paper\_tables.py
% Readability-oriented format, similar to the original manuscript style.
% Balanced accuracy is not included because it is unavailable in consolidated summary CSVs.

\begin{table*}[t]
\centering
\small
\caption{Performance comparison across datasets and tasks (10-fold CV means; readable manuscript format).}
\label{tab:readable-full-results}
\begin{tabular}{llccccc}
\toprule
\textbf{Dataset / Task} & \textbf{Model} & \textbf{Acc} & \textbf{F1} & \textbf{Prec} & \textbf{Rec} & \textbf{ROC-AUC} \\
\midrule
\multicolumn{7}{c}{\textbf{EVONS -- Fake-news Detection}} \\
\midrule
 & \textbf{mlp\_mistral} & 0.989 & 0.988 & 0.988 & 0.989 & 0.999 \\
 & mlp\_bert & 0.982 & 0.981 & 0.984 & 0.978 & 0.999 \\
 & logreg\_mistral & 0.975 & 0.973 & 0.978 & 0.967 & 0.997 \\
 & xgboost\_bert & 0.972 & 0.970 & 0.975 & 0.964 & 0.997 \\
 & logreg\_bert & 0.971 & 0.968 & 0.974 & 0.963 & 0.996 \\
 & xgboost\_mistral & 0.959 & 0.955 & 0.967 & 0.944 & 0.993 \\
 & rf\_bert & 0.949 & 0.944 & 0.953 & 0.935 & 0.989 \\
 & rf\_mistral & 0.931 & 0.922 & 0.963 & 0.885 & 0.984 \\
\midrule
\multicolumn{7}{c}{\textbf{EVONS -- Virality Prediction}} \\
\midrule
 & \textbf{gating\_mistral} & 0.950 & 0.312 & 0.525 & 0.229 & 0.881 \\
 & mlp\_avg\_eng & 0.951 & 0.087 & 0.631 & 0.048 & 0.868 \\
 & mlp\_source & 0.951 & 0.084 & 0.667 & 0.047 & 0.870 \\
 & gating\_bert & 0.950 & 0.006 & 0.063 & 0.003 & 0.868 \\
 & mlp\_text & 0.950 & 0.006 & 0.313 & 0.003 & 0.828 \\
\midrule
\multicolumn{7}{c}{\textbf{FakeNewsNet -- Fake-news Detection}} \\
\midrule
 & \textbf{rf\_bert} & 0.946 & 0.906 & 0.942 & 0.874 & 0.970 \\
 & mlp\_mistral & 0.945 & 0.903 & 0.959 & 0.857 & 0.981 \\
 & mlp\_bert & 0.940 & 0.895 & 0.941 & 0.857 & 0.975 \\
 & xgboost\_bert & 0.930 & 0.880 & 0.906 & 0.860 & 0.975 \\
 & rf\_mistral & 0.929 & 0.872 & 0.955 & 0.806 & 0.956 \\
 & logreg\_bert & 0.928 & 0.866 & 0.975 & 0.780 & 0.970 \\
 & xgboost\_mistral & 0.922 & 0.863 & 0.908 & 0.826 & 0.960 \\
 & logreg\_mistral & 0.869 & 0.721 & 1.000 & 0.566 & 0.982 \\
\midrule
\multicolumn{7}{c}{\textbf{FakeNewsNet -- Virality Prediction}} \\
\midrule
 & \textbf{xgboost\_bert} & 0.770 & 0.777 & 0.754 & 0.803 & 0.861 \\
 & mlp\_bert & 0.748 & 0.772 & 0.704 & 0.859 & 0.806 \\
 & logreg\_mistral & 0.755 & 0.772 & 0.723 & 0.828 & 0.808 \\
 & logreg\_bert & 0.749 & 0.766 & 0.720 & 0.819 & 0.809 \\
 & rf\_bert & 0.761 & 0.765 & 0.755 & 0.776 & 0.862 \\
 & mlp\_mistral & 0.750 & 0.760 & 0.735 & 0.797 & 0.809 \\
 & xgboost\_mistral & 0.736 & 0.740 & 0.731 & 0.750 & 0.831 \\
 & rf\_mistral & 0.749 & 0.740 & 0.768 & 0.716 & 0.849 \\
\bottomrule
\end{tabular}
\end{table*}


% Auto-generated by tools/generate_paper_tables.py

\begin{table*}[t]
\centering
\small
\caption{Top pairwise F1 comparisons (best model vs strongest alternatives per dataset/task).}
\label{tab:stats-top-pairs}
\begin{tabular}{lllcccl}
\toprule
Dataset/Task & Best & Compared to & $\Delta$F1 & 95\% CI & Holm-$p$ & Sig. \\
\midrule
evons/disinformation & mlp_mistral & rf_mistral & 0.0665 & [0.0648, 0.0679] & 0.0547 & no \\
evons/disinformation & mlp_mistral & rf_bert & 0.0445 & [0.0437, 0.0454] & 0.0547 & no \\
evons/disinformation & mlp_mistral & xgboost_mistral & 0.0334 & [0.0322, 0.0345] & 0.0547 & no \\
evons/virality & gating_mistral & mlp_text & 0.3060 & [0.2753, 0.3354] & 0.0195 & yes \\
evons/virality & gating_mistral & gating_bert & 0.3056 & [0.2786, 0.3340] & 0.0195 & yes \\
evons/virality & gating_mistral & mlp_source & 0.2279 & [0.1916, 0.2658] & 0.0195 & yes \\
fakenewsnet/disinformation & rf_bert & logreg_mistral & 0.1855 & [0.1597, 0.2137] & 0.0547 & no \\
fakenewsnet/disinformation & rf_bert & xgboost_mistral & 0.0431 & [0.0205, 0.0732] & 0.0820 & no \\
fakenewsnet/disinformation & rf_bert & logreg_bert & 0.0406 & [0.0184, 0.0638] & 0.1406 & no \\
fakenewsnet/virality & xgboost_bert & rf_mistral & 0.0374 & [0.0123, 0.0653] & 0.5156 & no \\
fakenewsnet/virality & xgboost_bert & xgboost_mistral & 0.0373 & [0.0162, 0.0573] & 0.3047 & no \\
fakenewsnet/virality & xgboost_bert & mlp_mistral & 0.0173 & [-0.0052, 0.0399] & 1.0000 & no \\
\bottomrule
\end{tabular}
\end{table*}



\paragraph{EVONS -- fake-news detection.}
Performance is near ceiling. The top mean F1 is achieved by \texttt{mlp\_mistral} (0.988), but corrected pairwise tests indicate no statistically reliable dominance over several strong alternatives.

\paragraph{EVONS -- virality prediction.}
This is the most difficult and most imbalanced condition. \texttt{gating\_mistral} reaches the highest F1 (0.312) and is the only setting where corrected significant separations are consistently observed against key alternatives (Table~\ref{tab:stats-top-pairs}).

\paragraph{FakeNewsNet -- fake-news detection.}
The best mean F1 is \texttt{rf\_bert} (0.906). Competing models remain close, and corrected significance is generally not observed among top contenders.

\paragraph{FakeNewsNet -- virality prediction.}
The best mean F1 is \texttt{xgboost\_bert} (0.777), with several alternatives within a narrow margin. Under correction, differences are not statistically significant.

\section{Discussion}
\label{sec:discussion}

Three conclusions emerge.

\begin{enumerate}
  \item \textbf{Task difficulty gap is stable.} Fake-news detection is consistently easier than virality prediction across datasets.
  \item \textbf{Label design strongly affects apparent performance.} EVONS q95 virality is substantially harder than FakeNewsNet median-split virality.
  \item \textbf{Statistical controls temper leaderboard interpretations.} In most settings, rank differences among top models are not robust after multiplicity correction; EVONS virality is the main exception.
\end{enumerate}

Overall, these results support reporting practices where protocol transparency and inferential checks are first-class components of model evaluation.

\section*{Ethics and Limitations}
Virality labels are operational proxies (q95 on EVONS; median split on FakeNewsNet) and should not be interpreted as universal definitions. Data availability and platform constraints can affect reproducibility and cross-paper comparability. We therefore emphasize transparent preprocessing assumptions, explicit label definitions, and conservative claims grounded in corrected statistical tests.

\section{Bibliographical References}\label{sec:reference}
\bibliographystyle{lrec2026-natbib}
\bibliography{lrec2026-example}

\end{document}
